\section{Conclusion}
\label{sec:conclusion}

This semester moved the localization pipeline to AWS and reduced the work it does on every commit. Profiling showed that
model inference takes 99.3\% of the translation stage. A one-change-at-a-time study then showed that the common intuition
about batching, that less padding means more speed, does not hold here: padding fell from 30\% to 4\% for a 3\% gain,
while the beam size and idle cores were the real costs. With beam~2, parallel batches and an output length cap, the
production pipeline translated the full twelve-language KDE4 corpus, 120,648 translations, in 3.6 hours with no failures,
about 2.5\x{} faster than at the start of the semester.

Scale exposed problems that small tests could not: broken placeholders and characters the model cannot write. Placeholder
protection raised survival from 87.7\% to 99.88\%, and untranslatable strings are now kept in English and reported rather
than shipped broken.

The redesigned Translation Memory changed the economics of the pipeline. Because it persists, records the engine that
produced each entry, loads once per language and reuses near-identical strings, the cost of a run now follows the size of
the change rather than the size of the application: 1\%, 5\% and 10\% commits ran 25\x{}, 14\x{} and 8.4\x{} faster with
90--98\% fewer model tokens and unchanged quality. Removing the remaining per-run set-up brought a typical 5\% commit to
3.1 minutes end to end, from roughly 1.9 hours at the start of the semester.
